% SGA 6, Exposición 0: fragmento español alineado con la fuente.
% Alcance: sección 4 completa, antes de la sección 5.
% Fuente francesa: sga6_fr_r1/src/fr.tex, líneas 363--416 inclusive.
% SHA-256 de la fuente francesa completa:
% 0D2158B9FE3B466309A92165727A0B191422E05FBAA0DACD5806E8322C54BA1A.

\subsection*{4.}
Finalmente, debemos dar algunas indicaciones sobre la manera en que todavía
se puede dar sentido a (1.2) en ausencia de una hipótesis de existencia de
módulos invertibles amplios sobre $X$ e $Y$.

\subsubsection*{4.1}
Observemos primero que \uline{para un esquema $X$ (o, más generalmente, un
topos localmente anillado) cualquiera, parece indudable que el «~buen~»
invariante que debe sustituir al anillo $K(X)$ de [RRR] no es el anillo
$K^\bullet(X)$ considerado anteriormente, definido en términos de módulos
localmente libres sobre $X$, sino el grupo $K(\Parf(X))$ considerado en 3.1,
definido en términos de complejos perfectos sobre $X$, grupo que aquí
denotaremos por $K^\bullet(X)$ y que distinguiremos del grupo «~ingenuo~» de
3.1, que denotamos por $K^\bullet_{\mathrm{naif}}(X)$}. Este punto de vista
viene impuesto por la necesidad de definir un homomorfismo (3.1) de imagen
directa para un morfismo $f:X\to Y$ propio y de tor-dimensión finita, de
esquemas noetherianos, digamos (el caso no noetheriano exige propiedades de
finitud locales algo más fuertes sobre $f$; cf. Exp. III). Con la definición
de $K^\bullet(X)$ adoptada ahora, este homomorfismo $f_*$ se define, en
efecto, de manera trivial mediante la fórmula (3.2), donde ahora
$\underline F$ es un complejo perfecto cualquiera sobre $X$, lo que implica
que $Rf_*(\underline F)$ también es perfecto.

\subsubsection*{4.2}
Esta nueva definición de $K^\bullet(X)$ plantea, sin embargo, un nuevo
problema de álgebra homológica: el de definir una $\lambda$-estructura sobre
$K^\bullet(X)$ compatible con su estructura natural de anillo, procedente
del producto tensorial en la categoría derivada. (NB: si $L_\bullet$ y $L'$
son dos complejos perfectos, lo mismo ocurre con su producto tensorial total
$L_\bullet\otimes^{\mathbf L}L'_\bullet$). Más precisamente, habría que definir
operaciones «~potencias exteriores~» en la categoría derivada $D^-(X)$ que
transformen complejos perfectos en complejos perfectos. Esta cuestión aún no
se ha aclarado en la actualidad y, como se dijo en la Introducción,
constituye la laguna más grave del Seminario. Una vez resuelta esta cuestión,
y procediendo como en 2.3, se obtiene una filtración natural sobre
$K^\bullet(X)$, que permite formar el anillo graduado asociado
$\Gr^\bullet(X)$, el cual desempeñará el papel del anillo de Chow, y definir
una teoría de clases de Chern
\[
c_i:K^\bullet(X)\longrightarrow \Gr^\bullet(X).
\]

\subsubsection*{4.3}
Para dar sentido a la fórmula (1.2) para un morfismo propio localmente de
intersección completa $f:X\to Y$ de esquemas noetherianos, aún hay que
incluir en el enunciado de esta fórmula las inclusiones no triviales (3.4)
(donde $d$ es la dimensión relativa virtual), a fin de poder definir (3.3),
y definir, finalmente, la clase tangente relativa virtual
\[
T_f\in K^\bullet(X).
\]
Cuando $f$ admite una factorización (2.1) mediante una inmersión en un
esquema relativo liso, la definición de $T_f$ no presenta ninguna dificultad
nueva. Por otra parte, en este caso puede reformularse de una manera más
conforme con el espíritu de las «~categorías derivadas~», considerando el
conocido homomorfismo canónico
\begin{equation}\tag{4.1}
\underline N_{X/X'}\xrightarrow{d}\Omega^1_{X'/Y}\otimes_{\underline O_{X'}}\underline O_X
\end{equation}
inducido por la diferencial exterior
\[
d_{X'/X}:\underline J\subset\underline O_{X'}\longrightarrow\Omega^1_{X'/Y},
\]
y el complejo de cadenas
\begin{equation}\tag{4.2}
L^{X/Y}_\bullet\in\operatorname{Ob}D(X)
\end{equation}
definido gracias a (4.1) «~prolongando por ceros~», complejo igual en grado
$1$, resp. $0$, a la fuente, resp. al destino, del homomorfismo (4.1). La
fórmula (2.2) adopta entonces la forma más agradable
\begin{equation}\tag{4.3}
\check T_f=\cl^\bullet\bigl(L^{X/Y}_\bullet\bigr),
\end{equation}
fórmula que tiene sentido en $K^\bullet(X)$ cuando $f$ es un morfismo
localmente de intersección completa, puesto que, en este caso, el complejo
$L^{X/Y}_\bullet$ es manifiestamente perfecto. La afirmación de invariancia
de $T_f$ respecto de la factorización elegida de $f$ puede precisarse
entonces fácilmente como una afirmación de invariancia para el complejo
$L^{X/Y}_\bullet$: salvo isomorfismo canónico en la categoría derivada
$D(X)$, este complejo también es independiente de la factorización elegida
(Exp. VIII).

\subsubsection*{4.4}
Cuando $f:X\to Y$ es un morfismo cualquiera de esquemas, todavía se puede
construir un objeto canónico $L^{X/Y}_\bullet$ de la categoría derivada
$D(X)$, definido salvo isomorfismo único, que sobre cada abierto afín de $X$
por encima de un abierto afín de $Y$ sea isomorfo al complejo construido
según el principio expuesto en el párrafo anterior (aunque, en el caso
general, conviene sustituir «~liso~» por «~formalmente liso~», de modo que
toda álgebra $B$ sobre un anillo $A$ aparece como cociente de un álgebra
formalmente lisa, por ejemplo un álgebra de polinomios). Para tal definición
no es posible proceder por simple «~pegado~» a partir de las construcciones
afines, pues los objetos de la categoría derivada $D(X)$ son de naturaleza
esencialmente no pegable.

He aquí una construcción general de $L^{X/Y}_\bullet$, válida más
generalmente para todo morfismo $f:X\to Y$ de topoi (conmutativamente)
anillados: sea $B=\underline O_X$, $A=f^{-1}(\underline O_Y)$, de modo que
$A$ es un anillo sobre el espacio $X$ y $B$ una $A$-álgebra. Sea $T\to B$
un homomorfismo de haces de conjuntos que engendre $B$ como $A$-álgebra, es
decir, tal que el homomorfismo de $A$-álgebras
\[
\varphi:C=A[T]\longrightarrow B
\]
correspondiente sea un epimorfismo de haces de conjuntos, donde $A[T]$
designa la $A$-álgebra conmutativa libre engendrada por $T$ (es decir, el haz
asociado al prehaz $U\longmapsto A(U)[T(U)]$ = $A(U)$-álgebra conmutativa
libre engendrada por $T(U)$).
Sea $J$ el núcleo de $\varphi$, y consideremos, como en 4.3, el complejo de
cadenas de longitud $1$ definido por el homomorfismo natural
\[
J/J^2\xrightarrow{d}\Omega^1_{C/A}\otimes_C B.
\]
No es difícil ver que, salvo isomorfismo canónico en la categoría $D(X)$, el
complejo así construido es independiente de la elección de $T$ y de
$T\to B$, de modo que merece la denominación de \uline{complejo cotangente
relativo $L^{X/Y}_\bullet$ de $X$ sobre $Y$}. Se comprueba asimismo que,
para todo abierto afín de $X$ por encima de un abierto afín de $Y$, viene
dado (salvo isomorfismo canónico) por el procedimiento de 4.3.

Cuando $f$ es localmente de intersección completa, se ve además que se
obtiene incluso un complejo \underline{perfecto}, lo que permite definir
también un elemento $T_f$ en $K^\bullet(X)$ mediante la fórmula (4.3).
Señalemos que, también aquí, no habría sido posible, en general, definir
$T_f$ como un elemento de $K^\bullet_{\mathrm{naif}}(X)$, y nada permite
suponer que $T_f$ sea \underline{globalmente} isomorfo en $D(X)$ a un
complejo acotado cuyas componentes sean localmente libres de tipo finito.
Esto constituye, pues, una justificación adicional del punto de vista
defendido en 4.1 a favor del $K^\bullet(X)$ «~sofisticado~», con preferencia
sobre $K^\bullet_{\mathrm{naif}}(X)$.

\subsubsection*{4.5}
A reserva de la cuestión planteada en 4.2, hemos reunido, pues, todos los
elementos de estructura necesarios para dar sentido a (1.2) para un
morfismo propio localmente de intersección completa de esquemas cualesquiera
(la hipótesis noetheriana hecha implícitamente en las secciones anteriores
puede eliminarse, de hecho, sin dificultad esencial). Por desgracia, esto no
significa que, aun con esta reserva, la fórmula en cuestión esté establecida,
ni siquiera cuando $Y$ es el espectro de un cuerpo algebraicamente cerrado
de característica $p>0$ y $X$ es un esquema algebraico propio y liso de
dimensión $3$, digamos. Véanse los comentarios de Exp. XIV
n\textsuperscript{o}~2.

